% ============================================================
%  06_tongket.tex – Chương 6: Tổng kết
% ============================================================
\chapter{Tổng Kết}

\section{Kết Quả Triển Khai}

Dự án đã xây dựng được một nền tảng thương mại điện tử cho Mọng Fruitboxz với hai khu vực sử dụng chính: giao diện mua sắm cho khách hàng và giao diện quản trị cho cửa hàng. Ở phía khách hàng, hệ thống hỗ trợ duyệt catalog, xem chi tiết sản phẩm, tìm kiếm, chọn biến thể, quản lý giỏ hàng, áp dụng khuyến mãi, tự động báo phí giao hàng và đặt đơn. Khách hàng đã đăng nhập có thêm khả năng quản lý tài khoản và theo dõi lịch sử mua hàng.

Ở phía vận hành, hệ thống cung cấp các chức năng quản lý sản phẩm, danh mục, đơn hàng, khách hàng, khuyến mãi, tồn kho, nguyên liệu, banner, blog, media, FAQ, cấu hình website và người dùng nội bộ. Một đặc thù nghiệp vụ được thể hiện trong đề tài là mô hình tồn kho theo công thức nguyên liệu: mỗi biến thể sản phẩm có thể được cấu thành từ nhiều nguyên liệu khác nhau, nhờ đó quá trình đặt hàng phản ánh gần hơn nghiệp vụ bán trái cây cắt sẵn và hộp quà.

Về mặt tổ chức hệ thống, dự án đã tách rõ giao diện, backend, dữ liệu và các dịch vụ hỗ trợ. Backend Medusa được mở rộng bằng các mô-đun riêng cho nội dung, phân quyền, nguyên liệu, yêu cầu mua số lượng lớn và giao hàng. Frontend sử dụng React để xây dựng trải nghiệm mua hàng và khu vực quản trị trong cùng một mã nguồn, đồng thời kết nối với các dịch vụ tìm kiếm, lưu trữ ảnh và chatbot tư vấn.

\section{Đối Chiếu Kết Quả Theo Nhóm Chức Năng}

\begin{longtable}{@{}L{0.27\textwidth} L{0.25\textwidth} L{0.40\textwidth}@{}}
  \toprule
  \textbf{Nhóm chức năng} & \textbf{Hiện trạng triển khai} & \textbf{Minh chứng và giới hạn} \\
  \midrule
  \endhead
  Catalog và tìm kiếm & Đã triển khai luồng chính & Sản phẩm, biến thể, danh mục, ảnh và bộ lọc đã có trên giao diện; backend có cơ chế trả kết quả dự phòng khi dịch vụ tìm kiếm không khả dụng. \\
  Giỏ hàng và đặt hàng & Đã triển khai đặt đơn; thanh toán còn ở mức ghi nhận & Người dùng có thể thêm sản phẩm, chỉnh số lượng, nhập thông tin nhận hàng, áp dụng khuyến mãi và tạo đơn. Phí giao hàng được tính lại ở backend; thanh toán trực tuyến chưa được tích hợp. \\
  Giao hàng theo khu vực & Đã có cấu hình và báo phí tự động & Hệ thống có gợi ý địa chỉ, định vị, báo phí theo quận/khoảng cách và lưu thông tin giao hàng trong đơn; lịch giao và đối tác vận chuyển chưa được mô hình hóa sâu. \\
  Quản trị vận hành & Đã có các màn hình nghiệp vụ chính & Khu vực quản trị bao phủ catalog, đơn hàng, khách hàng, khuyến mãi, tồn kho, nội dung và phân quyền; thao tác hàng loạt và cách gọi trạng thái vẫn cần làm chặt hơn. \\
  Nội dung và chăm sóc khách hàng & Đã có CMS và kênh tư vấn cơ bản & Banner, blog, FAQ, liên hệ và chatbot đã được kết nối với dữ liệu dự án; chất lượng tư vấn phụ thuộc vào việc bổ sung FAQ và dữ liệu sản phẩm khi vận hành. \\
  Tồn kho theo nguyên liệu & Đã có mô hình BOM & Công thức BOM liên kết sản phẩm bán ra với lượng nguyên liệu thực tế; nếu dùng cho cửa hàng thật, cần xử lý kỹ hơn phần cạnh tranh tồn kho và nhập/xuất kho. \\
  Báo cáo vận hành & Có dashboard tổng quan & Dashboard đã hiển thị một số chỉ số doanh thu, chi phí, lợi nhuận và công nợ; phân tích theo thời gian, nhóm sản phẩm và hiệu quả khuyến mãi là hướng mở rộng. \\
  \bottomrule
\end{longtable}

\section{Những Điểm Cần Hoàn Thiện}

\begin{itemize}[leftmargin=2em, itemsep=5pt]
  \item \textbf{Thanh toán trực tuyến:} hiện đơn hàng chủ yếu dừng ở ghi nhận đơn và hướng dẫn chuyển khoản. Nếu triển khai thực tế, cần kết nối cổng thanh toán phù hợp và có cơ chế xác nhận trạng thái thanh toán.
  \item \textbf{Thông báo sau đặt hàng:} email xác nhận và thông báo cho nhân viên cửa hàng cần được kết nối với nhà cung cấp gửi thư hoặc kênh nhắn tin thật.
  \item \textbf{Quy trình giao vận:} trạng thái giao hàng đã được mô hình hóa ở mức cơ bản; có thể bổ sung lịch giao, khu vực phục vụ, người phụ trách và ghi chú xử lý đơn.
  \item \textbf{Báo cáo kinh doanh:} dashboard hiện cho cái nhìn tổng quan; phần phân tích doanh thu, chi phí, lợi nhuận, tồn kho và hiệu quả khuyến mãi theo thời gian còn có thể trình bày sâu hơn.
  \item \textbf{Trải nghiệm quản trị:} một số màn hình quản trị có nhiều chức năng; nên thống nhất lại cách lọc, tìm kiếm, phân trang và thao tác nhanh để nhân viên sử dụng thuận tiện hơn.
  \item \textbf{Nội dung tư vấn:} chatbot đã hỗ trợ hỏi đáp cơ bản; chất lượng câu trả lời sẽ tốt hơn khi FAQ, dữ liệu sản phẩm và kịch bản tư vấn được bổ sung thường xuyên từ quá trình vận hành thực tế.
\end{itemize}

\section{Hướng Phát Triển}

\begin{itemize}[leftmargin=2em, itemsep=5pt]
  \item Tích hợp thanh toán trực tuyến, đối soát giao dịch và hiển thị trạng thái thanh toán rõ ràng cho khách hàng.
  \item Hoàn thiện thông báo đơn hàng qua email, SMS hoặc kênh nội bộ để nhân viên xử lý đơn nhanh hơn.
  \item Mở rộng cấu hình giao hàng theo khung giờ, khu vực phục vụ, phí theo thời điểm và lựa chọn đơn vị vận chuyển.
  \item Phát triển báo cáo quản trị theo ngày, tuần, tháng, nhóm sản phẩm, khách hàng và hiệu quả chương trình khuyến mãi.
  \item Cải thiện chatbot theo hướng ghi nhận câu hỏi phổ biến, gợi ý sản phẩm chính xác hơn và hỗ trợ nhân viên bổ sung FAQ từ trang quản trị.
  \item Cải thiện phiên bản mobile và nội dung SEO để website dùng tốt hơn trong bối cảnh bán hàng thật.
\end{itemize}

\section{Kiến Thức và Kinh Nghiệm Thu Được}

Quá trình thực hiện giúp sinh viên củng cố cách xây dựng một hệ thống thương mại điện tử theo kiến trúc headless commerce, cách tách giao diện khách hàng khỏi backend, cách tổ chức dữ liệu sản phẩm--đơn hàng--tồn kho và cách mở rộng Medusa bằng các mô-đun riêng. Đề tài cũng là dịp luyện kỹ năng đọc mã nguồn lớn, đối chiếu giữa giao diện và nghiệp vụ, vẽ sơ đồ phân tích--thiết kế, cũng như trình bày kết quả thành một báo cáo có cấu trúc.

Một kinh nghiệm quan trọng là chức năng bán hàng không chỉ nằm ở trang sản phẩm hay nút đặt mua. Để tạo được một luồng mua hàng hoàn chỉnh, hệ thống phải phối hợp nhiều phần: dữ liệu sản phẩm, giá bán, tồn kho, phí giao hàng, khuyến mãi, tài khoản khách hàng, trạng thái đơn và công cụ cho nhân viên vận hành.

\section{Kết Luận}

Mọng Fruitboxz đã triển khai được các thành phần chính của một website bán hàng cho cả khách hàng và cửa hàng. Hệ thống thể hiện được đặc thù của sản phẩm trái cây qua biến thể, hộp quà, giao hàng theo khu vực, chatbot tư vấn và tồn kho theo nguyên liệu. Dù vẫn còn các phần cần hoàn thiện như thanh toán trực tuyến, thông báo sau đặt hàng và báo cáo chuyên sâu, phiên bản hiện tại vẫn đủ để làm nền cho các bước phát triển tiếp theo của hệ thống.

\section{Tài Liệu Tham Khảo}

\begin{enumerate}[leftmargin=2em, itemsep=4pt]
  \item Medusa Documentation. \url{https://docs.medusajs.com/}
  \item React Documentation. \url{https://react.dev/}
  \item Vite Documentation. \url{https://vite.dev/guide/}
  \item Tailwind CSS Documentation. \url{https://tailwindcss.com/docs}
  \item PostgreSQL 15 Documentation. \url{https://www.postgresql.org/docs/15/}
  \item Redis Documentation. \url{https://redis.io/docs/latest/}
  \item Meilisearch Documentation. \url{https://www.meilisearch.com/docs/}
  \item MinIO Documentation. \url{https://min.io/docs/minio/}
  \item Docker Compose Documentation. \url{https://docs.docker.com/compose/}
  \item Groq API Documentation. \url{https://console.groq.com/docs/}
  \item Mã nguồn Mọng Fruitboxz. \url{https://github.com/pathust/mong.fruitboxz}
\end{enumerate}
